%%%%%%%%%%%%%%%%%%%%%%%%%%%%%%%%%%%%%%%%%%%%%%%%%%%%%%%
%%%% DEFINE CONFERENCE SPECIFIC PACKAGES AND STYLE %%%%
%%%%%%%%%%%%%%%%%%%%%%%%%%%%%%%%%%%%%%%%%%%%%%%%%%%%%%%

% some packages that might cause errors
% NOTE: hyperref is loaded by automl.sty; add it here only when not using automl
\usepackage[utf8]{inputenc}
\usepackage{url}

% math
\usepackage{amsmath}
\let\Bbbk\relax  % newtxmath (via automl.sty) defines \Bbbk before amssymb does
\usepackage{amssymb}
\usepackage{bbm}
\usepackage{bm}
\usepackage{cancel}
\usepackage{mathtools}

% algorithm
\usepackage{ascmac}  % ascmac must come before `algorithm' and `algpseudocode'
\usepackage{algorithm}
\usepackage{algpseudocode}

% fancy viz
% NOTE: fancyhdr is loaded and configured by automl.sty
\usepackage{tcolorbox}
\usepackage{tikz}

% utils
\usepackage{booktabs}
\usepackage{comment}
\usepackage{lastpage}
\usepackage{listings}
\usepackage{multibib}
\usepackage{multirow}
\usepackage[super]{nth}
% NOTE: subfig conflicts with subcaption loaded by automl.sty; use \subcaptionbox instead

\tcbuselibrary{breakable, skins, theorems}

\allowdisplaybreaks

% a mathematical statement that is proved using rigorous mathematical reasoning.
% In a mathematical paper, the term theorem is often reserved for the most important results.
\newtheorem{theorem}{Theorem}

% a result in which the (usually short) proof relies heavily on a given theorem.
\newtheorem{corollary}{Corollary}

% a proved and often interesting result, but generally less important than a theorem.
\newtheorem{proposition}{Proposition}

% a minor result whose sole purpose is to help in proving a theorem.
% It is a stepping stone on the path to proving a theorem.
% Very occasionally lemmas can take on a life of their own.
\newtheorem{lemma}{Lemma}

% an assertion that is then proved.
% It is often used like an informal lemma.
\newtheorem{claim}{Claim}

% a precise and unambiguous description of the meaning of a mathematical term.
% It characterizes the meaning of a word by giving all the properties and only those properties that must be true.
\newtheorem{definition}{Definition}

% a statement that is assumed to be true without proof.
% These are the basic building blocks from which all theorems are proved.
\newtheorem{assumption}{Assumption}

\newtheorem{proof}{Proof}

\def\qed{\hfill $\Box$}

% remove the end from algorithm
\algtext*{EndFor}
\algtext*{EndWhile}
\algtext*{EndIf}
\algtext*{EndProcedure}
\algtext*{EndFunction}
\algnewcommand{\LineComment}[1]{\State \(\triangleright\) #1}
\newcommand{\Break}{\textbf{break}}
\newcommand{\Continue}{\textbf{continue}}

% Add `References` for Appendix
\newcites{appx}{References}

%%%%%%%%%%%%%%%%%%%%%%%%%%%
%%%% PAPER STATUS FLAGS %%%%
%%%%%%%%%%%%%%%%%%%%%%%%%%%

\newif\ifunderreview
\newif\ifsubmission
\newif\ifappendix

% Whether under review or not
% \underreviewtrue
\underreviewfalse

% Whether submission version or not
% \submissiontrue
\submissionfalse

% Whether to include appendix or not
\appendixtrue
% \appendixfalse

\ifsubmission
\newcommand{\todo}[1]{}
\newcommand{\replace}[2]{}
\newcommand{\sw}[1]{}

\else
% ToDo
\newcommand{\todo}[1]{\textbf{\textcolor{red}{[TODO: #1]}}}

% deletion
\newcommand{\replace}[2]{\textbf{\textcolor{red}{[del: \cancel{#1}]}}\textbf{\textcolor{blue}{[new: #2]}}}

% Shuhei comment
\newcommand{\sw}[1]{\textbf{\textcolor{cyan}{[SW: #1]}}}

% show label names inline
\usepackage[inline]{showlabels}

\fi

%%%%%%%%%%%%%%%%%%%%%%%%%%%
%%%% USER COMMANDS      %%%%
%%%%%%%%%%%%%%%%%%%%%%%%%%%

% The labeling for the section numbers in appendix
\makeatletter
\newcommand{\customlabel}[2]{%
   \protected@write \@auxout {}{\string \newlabel {#1}{{#2}{\thepage}{#2}{#1}{}} }%
   \hypertarget{#1}{}
}
\makeatother

\newcommand{\citewithname}[2]{#1 \textit{et al.} \shortcite{#2}}

% argmin/argmax
\newcommand{\argmin}{\mathop{\mathrm{argmin}}}
\newcommand{\argmax}{\mathop{\mathrm{argmax}}}
\newcommand{\indic}[1]{\mathbb{I}[#1]}
\newcommand{\comb}[2]{{}_#1 \mathrm{C} _#2}
\newcommand{\ceil}[1]{\lceil #1 \rceil}
\newcommand{\floor}[1]{\lfloor #1 \rfloor}

\newcommand{\blue}[1]{\textcolor{blue}{#1}}
\newcommand{\brown}[1]{\textcolor{brown}{#1}}
\newcommand{\cyan}[1]{\textcolor{cyan}{#1}}
\newcommand{\darkgray}[1]{\textcolor{darkgray}{#1}}
\newcommand{\gray}[1]{\textcolor{gray}{#1}}
\newcommand{\green}[1]{\textcolor{green}{#1}}
\newcommand{\lightgray}[1]{\textcolor{lightgray}{#1}}
\newcommand{\lime}[1]{\textcolor{lime}{#1}}
\newcommand{\magenta}[1]{\textcolor{magenta}{#1}}
\newcommand{\olive}[1]{\textcolor{olive}{#1}}
\newcommand{\orange}[1]{\textcolor{orange}{#1}}
\newcommand{\pink}[1]{\textcolor{pink}{#1}}
\newcommand{\purple}[1]{\textcolor{purple}{#1}}
\newcommand{\red}[1]{\textcolor{red}{#1}}
\newcommand{\teal}[1]{\textcolor{teal}{#1}}
\newcommand{\violet}[1]{\textcolor{violet}{#1}}
\newcommand{\white}[1]{\textcolor{white}{#1}}
\newcommand{\yellow}[1]{\textcolor{yellow}{#1}}

\newcommand{\secref}[1]{Section~\ref{#1}}
\renewcommand{\eqref}[1]{Eq.~(\ref{#1})}
\newcommand{\figref}[1]{Figure~\ref{#1}}
\newcommand{\tabref}[1]{Table~\ref{#1}}

\DeclareFixedFont{\ttb}{T1}{txtt}{bx}{n}{12} % for bold
\DeclareFixedFont{\ttm}{T1}{txtt}{m}{n}{12}  % for normal

\definecolor{deepblue}{rgb}{0,0,0.5}
\definecolor{deepred}{rgb}{0.6,0,0}
\definecolor{deepgreen}{rgb}{0,0.5,0}

\lstset{
  language=Python,
  basicstyle=\ttm,
  morekeywords={self},              % Add keywords here
  keywordstyle=\ttb\color{deepblue},
  emph={MyClass,__init__},          % Custom highlighting
  emphstyle=\ttb\color{deepred},    % Custom highlighting style
  stringstyle=\color{deepgreen},
  frame=tb,                         % Any extra options here
  showstringspaces=false,
  commentstyle=\color{deepgreen}
}

% \author[1,$\ast$]{\nameemail{Author 1}{email1@example.com}}
% \author[1,2,$\ast$]{\nameemail{Author 2}{email2@example.com}}
% \author[2]{\nameemail{Author 3}{email3@example.com}}

% \affil[$\ast$]{Equal contribution.}
% \affil[1]{Institution 1}
% \affil[2]{Institution 2}

\title{title}

% define PDF metadata, aids in accessibility of the resulting PDF
\ifunderreview
% no metadata for anonymity.
\else
\hypersetup{%
  pdfauthor={AutoML},
  pdftitle={Documentation for the automl Package},
  pdfsubject={Documentation for automl Package},
  pdfkeywords={AutoML, LaTeX, style}
}
\fi
